\documentclass{article}
\usepackage{amsmath,amssymb}
\begin{document}
\section*{Common train timetable with scenario-dependent integer passenger flows}
The original extraction is retained separately. This complete event-indexed
form specifies its timetable, transfer and passenger constraints and both
objective terms. Source text on printed page 73 explicitly rounds alighting
up and transfers down; the old extraction omitted that convention.

Lines are $l\in L$, directional station visits $i=1,\ldots,I_l$ and demand
scenarios $s\in S$, of probability $\pi_s$ with $\sum_s\pi_s=1$.
The time step is $\tau$ minutes and the normalized horizon is $t=1,\ldots,T$.
Input station offsets $e_{li}$, dwell times $d_{li}$, headways
$h_l^{\min},h_l^{\max}$ and transfer walks $f_{li,l'j}$ are in time-step
units. Each line has capacity $C_l$, service count bounds
$K_l^{\min},K_l^{\max}$ and service cost $c_l$.
$p_{slit}$ is external arrival count, $\alpha_{sli}\in[0,1]$ the
alighting fraction and $\rho_{li,l'j}\in[0,1]$ the transfer fraction,
with outgoing transfer fractions summing to at most one.
Time-constant fractions are a property of the supplied instances.
The last visit has $\alpha=1$ and zero external arrivals.

\subsection*{Timetable and transfer events}
Choose common integer service count $K_l$ and departure indices
$t_{lk}\in\{1,\ldots,T\}$ for $k=1,\ldots,K_l$, in increasing order.
Equivalently $x_{lkt}=1$ exactly when $t=t_{lk}$; inactive service indices
have all-zero $x$.
\begin{align}
K_l^{\min}\le K_l&\le K_l^{\max},\\
h_l^{\min}\le t_{l1}&\le h_l^{\max},\\
h_l^{\min}\le t_{lk}-t_{l,k-1}&\le h_l^{\max}\quad(k>1),\\
\sum_kx_{lkt}\le1,\quad \sum_tx_{lkt}\le1,\quad
\sum_tx_{lkt}&\le\sum_tx_{l,k-1,t}\quad(k>1).
\end{align}
The first-time convention is $T_0=1$, so the source first-time bounds
$T_0+h-1$ become $h$. Boarding at visit $i$ occurs at absolute time
$t_{lk}+e_{li}$; alighting occurs at $t_{lk}+e_{li}-d_{li}$.
For every permitted transfer pair $(li,l'j)$, its event for train $k$ is
at absolute time $t_{lk}+e_{li}-d_{li}+f_{li,l'j}$, or receiving-station
normalized time
\[
t'=t_{lk}+e_{li}-d_{li}+f_{li,l'j}-e_{l'j}.
\]
Define $y_{li,l'j,k,t'}=1$ exactly at that event, and zero otherwise.
Thus each active train has one event per permitted pair and none for
unconnected pairs; $\sum_k y_{li,l'j,k,t'}\le1$ follows from the positive
headway and a fixed transfer offset. All scenario decisions use these
same $x,y$.

\subsection*{Complete passenger constraints}
The model includes all external-arrival, alighting, transfer-arrival,
boarding and sampling events. Queues initially equal zero.
At equal absolute time the order is alight, external arrival, transfer
arrival, board, then sample. Each train has initial load $u_{slk0}=0$.
Let $u^-_{slki}=u_{slk,i-1}$, alighting count $a_{slki}$, boarding count
$b_{slki}$ and load after departure $u_{slki}$. For every active train:
\begin{align}
a_{slki}&=\lceil\alpha_{sli}u^-_{slki}\rceil,\qquad
0\le a_{slki}\le u^-_{slki},\\
v_{s,li,l'j,k}&=\lfloor\rho_{li,l'j}a_{slki}\rfloor,\\
u_{slki}&=u^-_{slki}-a_{slki}+b_{slki},\qquad0\le u_{slki}\le C_l,\\
0\le b_{slki}&\le Q^-_{sli,\ t_{lk}+e_{li}},\\
Q^+_{sli,\ t_{lk}+e_{li}}&=Q^-_{sli,\ t_{lk}+e_{li}}-b_{slki}.
\end{align}
At each external event $Q^+=Q^-+p_{slit}$; at each transfer-arrival
event $Q^+=Q^-+v_{s,li,l'j,k}$ at its receiving queue.
Queues remain unchanged at other events and always nonnegative.
All queue, load, boarding, alighting and transfer counts are integers.
Inactive services have zero boarding/load/transfer variables.
Alighting passengers not assigned a transfer leave the system; transfers
remain in transit until their specified receiving event. Passengers present at the horizon are carried into the terminal state. Capacity holds before and after every
boarding, so $b\le C_l$ and $v\le C_l y$ follow as well.
For any event, the total external arrivals so far equal exited passengers
plus all waiting queues, onboard loads and passengers still transferring.
The terminal alighting rule and zero terminal arrivals make final onboard
loads zero.

\subsection*{Objective and unequal normalized horizons}
Let $w_{slit}$ be the queue sampled after service at absolute time
$t+e_{li}$. The original objective counts external arrivals with a
half-step waiting approximation and waiting queues at normalized samples.
Some transfer events can fall beyond the receiving visit's last sample.
Define $W_s^{\rm end}=\sum_{li}w_{sliT}$ and
$W_s^{\rm late}=\sum_{li}Q_{sli}^{\rm eventual}-W_s^{\rm end}\ge0$,
where eventual queues are after every scheduled train and transfer event.
No train boards after its own visit's final sample; hence this difference
contains exactly the additional late transfers. They receive the same
terminal waiting penalty.
\begin{align}
Z_{\rm TSC}&=\sum_l c_lK_l,\\
PWT_s&=\frac{\tau}{2}\sum_{li,t}p_{slit}
+\tau\sum_{li,t<T}w_{slit}
+\epsilon(W_s^{\rm end}+W_s^{\rm late}),\\
\min Z&=Z_{\rm TSC}+\mu\sum_s\pi_sPWT_s .
\end{align}
The late-transfer term is an explicit boundary completion of the
extracted model, using the explicit transfer and terminal-state definitions above.

\subsection*{Algorithm and input scope}
The code searches service counts and departure-time neighborhoods, and
constructs feasible passenger flows by boarding as many as available
capacity allows. Greedy boarding and the search neighborhoods are algorithm
choices for candidate construction. The full supplied
scenario set is evaluated for every candidate. This is a verified feasible
heuristic with full scenario evaluation. The two delivered
instances retain respectively one line/five demand scenarios and two
lines/one scenario with 44 physical stations in total. Their time horizons
are the supplied input horizons, specified separately for each instance.
\end{document}
